% Part: proof-theory
% Chapter: cut-elimination
% Section: ce-topmost

\documentclass[../../../include/open-logic-section]{subfiles}

\begin{document}

\olfileid[id]{pt}{cut}{top}

\olsection{Menghapus Potongan Teratas}

\begin{defn}
Perhatikan !!a{proof}~$\pi$ yang berakhir dengan aturan \CutCS{},
\[
\AxiomC{}
\RightLabel{$\pi_1$}
\Deduce$\Gamma \fCenter \Delta, !A$
\AxiomC{}
\RightLabel{$\pi_2$}
\Deduce$!A, \Gamma \fCenter \Delta$
\RightLabel{\CutCS}
\BinaryInf$\Gamma \fCenter \Delta$
\DisplayProof
\]
dan tidak memuat inferensi \CutCS{} lain. \emph{Tinggi potongan} dari~$\pi$
adalah $\cheight{\pi} = \pheight{\pi_1} + \pheight{\pi_2}$.
\emph{Peringkat potongan} dari~$\pi$ adalah $\cutrank{\pi} = \depth{!A}$.
\end{defn}

\begin{lem}\ollabel{lem:cut-adm-G3c}
Aturan \CutCS{} dapat diterima dalam~\Log{G3c}. Dengan kata lain, jika
$\pi$ adalah !!a{proof} yang berakhir dengan satu~\CutCS{} dan selebihnya
hanya menggunakan aturan-aturan~$\Log{G3c}$, maka terdapat !!a{proof}~$\pi'$
atas sekuen akhir yang sama dalam~\Log{G3c}.
\end{lem}

\begin{proof}
Bukti berjalan melalui induksi ganda pada !!{height} dan peringkat potongan.
!!{proof}~$\pi$ berakhir dengan:
\[
\AxiomC{}
\RightLabel{$\pi_1$}
\Deduce$\Gamma \fCenter \Delta, !A$
\AxiomC{}
\RightLabel{$\pi_2$}
\Deduce$!A, \Gamma \fCenter \Delta$
\RightLabel{\CutCS}
\BinaryInf$\Gamma \fCenter \Delta$
\DisplayProof
\]

Dasar induksi: $\cheight{\pi} = 0$ dan $\cutrank{\pi} = 0$. Karena
$\cheight{\pi} = \pheight{\pi_1} + \pheight{\pi_2} = 0$, baik $\Gamma
\Sequent \Delta, !A$ maupun $!A, \Gamma \Sequent \Delta$ merupakan
aksioma. Ada dua kasus: (a)~$!A$ tidak prinsipal dalam salah satunya, atau
(b)~$!A$ prinsipal dalam keduanya. Di bawah ini kita akan menunjukkan bahwa
dalam kedua kasus $\Gamma \Sequent \Delta$ harus merupakan aksioma (dan
kita tidak memerlukan hipotesis induksi untuk menunjukkannya).

Sekarang kita membedakan sejumlah kemungkinan kasus berdasarkan apakah
$\pi_1$ dan $\pi_2$ merupakan aksioma atau berakhir dengan inferensi, dan
apakah $!A$ prinsipal dalam inferensi itu. Kasus A dan~B mencakup dasar
induksi. Dalam langkah induksi, kita mengasumsikan $\cheight{\pi} > 0$ atau
$\cutrank{\pi} > 0$ (atau keduanya). Kita dapat menggunakan hipotesis
induksi, yakni bahwa \CutCS{} terakhir dapat dihapus dari setiap !!{proof}
yang berakhir dengan satu~\CutCS{} yang mempunyai peringkat potongan
$= \cutrank{\pi}$ dan !!{height} potongan $< \cheight{\pi}$, atau mempunyai
peringkat potongan~$< \cutrank{\pi}$. Kasus $\cheight{\pi} = 0$ (kedua
premis merupakan aksioma) dicakup oleh kasus~A dan~B. Kasus
$\cheight{\pi} > 0$ dicakup oleh kasus C--E.

\paragraph{Kasus A:}
Suatu $!B$ muncul baik dalam $\Gamma$ maupun~$\Delta$. Maka $\Gamma
\Sequent \Delta$ merupakan aksioma (jika $!B$ atomik), atau mempunyai
!!a{proof}~$\pi'$ dalam \Log{G3c} tanpa \CutCS{} menurut
\olref[seq][ex]{prop:G3c-axioms}. Ini mencakup kasus~(a) dari dasar induksi,
sebab jika $!A$ tidak prinsipal dalam $\Gamma \Sequent \Delta, !A$ atau
tidak prinsipal dalam $!A, \Gamma \Sequent \Delta$, dan sekuen-sekuen itu
merupakan aksioma, maka suatu $!B$ atomik lain harus muncul dalam $\Gamma$
maupun~$\Delta$. Ini juga mencakup langkah induksi jika salah satu premis
merupakan aksioma tempat $!A$ tidak prinsipal (bahkan jika premis lain
merupakan kesimpulan suatu aturan).

\paragraph{Kasus B:} Kedua premis merupakan aksioma, dan $!A$ prinsipal
serta karenanya atomik. Perhatikan premis kiri, $\Gamma \Sequent \Delta,
!A$. Karena $!A$ prinsipal, formula itu harus muncul dalam $\Gamma$ (jika
tidak, $\Gamma \Sequent \Delta, !A$ bukan aksioma). Demikian pula, $!A$
harus muncul dalam $\Delta$; jika tidak, $!A, \Gamma \Sequent \Delta$ bukan
aksioma. Jadi, $!A$ atomik dan muncul baik dalam $\Gamma$ maupun $\Delta$,
yakni $\Gamma \Sequent \Delta$ merupakan aksioma. Ini mencakup kasus~(b)
dari dasar induksi.

\begin{center}
Kasus alternatif A dan B
\end{center}

\paragraph{Kasus A:} Salah satu premis potongan merupakan aksioma berbentuk
$!C, \Gamma' \Sequent \Delta', !C$ dengan $!C$ atomik. Jika $!C$ bukan
!!{formula} potongan~$!A$, maka $!C$ harus muncul baik dalam $\Gamma$ maupun
$\Delta$. Dalam hal ini, $\Gamma \Sequent \Delta$ merupakan aksioma dan
dapat kita ambil sebagai~$\pi'$. Jika $!C \ident !A$ adalah !!{formula}
potongan, maka (jika premis kiri merupakan aksioma) $!A \in \Gamma$ dan
$\Gamma = \Gamma', !A$, atau (jika premis kanan merupakan aksioma) $!A \in
\Delta$ dan $\Delta = \Delta', !A$. Dalam kasus pertama, $\pi_2$ merupakan
!!a{proof} dari $!A, !A, \Gamma' \Sequent \Delta$, dan menurut dapat
diterimanya \LeftR{\Contraction}
(\olref[seq][inv]{prop:G3c-cont-adm}) terdapat !!a{proof} dari $\Gamma
\Sequent \Delta$. Dalam kasus kedua, $\pi_1$ merupakan !!a{proof} dari
$\Gamma \Sequent \Delta', !A, !A$, dan menurut dapat diterimanya
\RightR{\Contraction} kita memperoleh !!a{proof} dari $\Gamma \Sequent
\Delta$.

\paragraph{Kasus B:} Salah satu premis merupakan aksioma berbentuk
$\lfalse, \Gamma' \Sequent \Delta'$. Jika premis kiri merupakan aksioma
semacam itu, maka $\lfalse \in \Gamma$ dan $\Gamma \Sequent \Delta$ juga
merupakan aksioma. Jika premis kanan merupakan aksioma semacam itu, maka
$\lfalse \in \Gamma$ (dalam hal ini $\Gamma \Sequent \Delta$ kembali
merupakan aksioma), atau $!A \ident \lfalse$. Dalam kasus terakhir,
$\pi_1$ merupakan !!a{proof} dari $\Gamma \Sequent \Delta, \lfalse$.
Kita dapat mengubahnya menjadi !!a{proof} dari~$\Gamma \Sequent \Delta$
dengan menghapus satu salinan~$\lfalse$ dari konsekuen setiap sekuen
dalam~$\pi_1$. (Latihan: periksa hal ini melalui induksi pada
$\pheight{\pi_1}$.)

Sekarang andaikan baik kasus~A maupun kasus~B tidak berlaku. Dalam
kasus-kasus yang tersisa, $\cheight{\pi} > 0$, yakni setidaknya satu premis
merupakan kesimpulan suatu inferensi.

\paragraph{Kasus C dan D: Memutasikan potongan.}
Kemungkinan-kemungkinan dalam kasus~C dan~D semuanya mempunyai pola umum.
Kita mulai dengan !!a{proof} yang berakhir dengan~\CutCS{}, dengan salah
satu premis diperoleh melalui aturan~$R$ tempat !!{formula} potongan~$!A$
tidak prinsipal (suatu !!{formula} lain~$!D$ yang prinsipal). Dalam
!!{proof} ini, \CutCS{} mengikuti aturan~$R$, dan secara skematis tampak
sebagai berikut:
\[
\AxiomC{}
\RightLabel{$\pi_1$}
\Deduce$\Gamma \fCenter \Delta', !D, !A$
\AxiomC{}
\RightLabel{$\pi_3$}
\Deduce$!A, \Gamma, \Pi \fCenter \Delta', \Lambda$
\RightLabel{$R$}
\UnaryInf$!A, \Gamma \fCenter \Delta', !D$
\RightSubproofLabel{$\pi_2$}
\RightLabel{\CutCS}
\BinaryInf$\Gamma \fCenter \Delta', !D$
\DisplayProof
\]
Ini hanya mencakup kasus ketika $R$ merupakan aturan kanan dengan satu
premis, $A$ bukan !!{formula} prinsipal dari~$R$, dan !!{formula}~$!D$ yang
\emph{prinsipal} muncul dalam sukseden premis kanan dari~\CutCS. Jika $!D$
muncul dalam anteseden (yakni $R$ merupakan aturan kiri) atau dalam premis
kiri~\CutCS, gambarnya tentu berbeda. !!{formula}s dalam $\Pi$ dan~$\Lambda$
mewakili !!{formula}s aktif aturan~$R$.

Kita membalik urutan ini dan memperhatikan !!a{proof} tempat $R$
mengikuti~\CutCS{}:
\[
\AxiomC{}
\RightLabel{$\pi_1'$}
\Deduce$\Gamma, \Pi \fCenter \Delta', !D, \Lambda, !A$
\AxiomC{}
\RightLabel{$\pi_3'$}
\Deduce$!A, \Gamma, \Pi \fCenter \Delta', !D, \Lambda$
\RightLabel{\CutCS}
\BinaryInf$\Gamma, \Pi \fCenter \Delta', !D, \Lambda$
\RightLabel{$R$}
\UnaryInf$\Gamma \fCenter \Delta', !D, !D$
\DisplayProof
\]
Kita menggunakan pelemahan yang mempertahankan !!{height} untuk memperoleh
$\pi_1'$ dan $\pi_3'$ masing-masing dari $\pi_1$ dan $\pi_3$.

Karena kita menukar urutan \CutCS{} dan~$R$, langkah ini disebut
``memutasikan suatu potongan ke atas melewati~$R$.'' Langkah tersebut
mengurangi tinggi !!{proof} yang berakhir pada premis kanan~\CutCS, dan
karena itu mengurangi !!{height} potongan. Artinya, sub-!!{proof}s yang
berakhir pada \CutCS{} baru dapat diasumsikan mempunyai !!{height} tidak
lebih besar daripada $\pi_1$ dan $\pi_3$, sehingga !!{height} potongan
berkurang (kita telah menghapus satu inferensi di sebelah kanan). Kini
hipotesis induksi berlaku dan memberi tahu kita bahwa inferensi \CutCS{}
dapat dihapus. Terakhir, kita menghapus salinan tambahan~$!D$ menggunakan
dapat diterimanya~$\RightR{\Contraction}$.

\paragraph{Kasus C:}
Setidaknya satu dari $\pi_1$ dan $\pi_2$ berakhir dengan aturan satu-premis
$R$ tempat $!A$ tidak prinsipal. Terdapat total 18 kasus, bergantung pada
apakah $!A$ tidak prinsipal dalam premis kiri atau kanan dari~\CutCS, dan
pada aturan mana dari sembilan aturan satu-premis (\LeftR{\lnot},
\RightR{\lnot}, \RightR{\lor}, \LeftR{\land}, \RightR{\lif}, dan keempat
aturan kuantor) yang merupakan aturan~$R$. Kita hanya membahas dua contoh:
ketika $!A$ tidak prinsipal dalam inferensi terakhir~$\pi_2$, yang berakhir
dengan~$\RightR{\lif}$ atau dengan~\LeftR{\lexists}. Kasus-kasus lain analog
dan diserahkan sebagai latihan.

Andaikan $\pi$ berakhir dengan:
\[
\AxiomC{}
\RightLabel{$\pi_1$}
\Deduce$\Gamma \fCenter \Delta', !B \lif !C, !A$
\AxiomC{}
\RightLabel{$\pi_3$}
\Deduce$!A, !B, \Gamma \fCenter \Delta', !C$
\RightLabel{$\RightR{\lif}$}
\UnaryInf$!A, \Gamma \fCenter \Delta', !B \lif !C$
\RightSubproofLabel{$\pi_2$}
\RightLabel{\CutCS}
\BinaryInf$\Gamma \fCenter \Delta', !B \lif !C$
\DisplayProof
\]
dengan $\Delta = \Delta', !B \lif !C$. Kita menerapkan pelemahan yang
mempertahankan !!{height} (\olref[seq][adm]{prop:weak-G3c-adm}) pada $\pi_1$
untuk memperoleh !!a{proof} $\pi_1'$ dari $!B, \Gamma \fCenter \Delta',
!B \lif !C, !C, !A$ (kita melemahkan dengan $!B$ di kiri dan $!C$ di
kanan). Demikian pula, kita menerapkan
\olref[seq][adm]{prop:weak-G3c-adm} pada~$\pi_3$ untuk memperoleh
!!a{proof}~$\pi_3'$ dari $!A, !B, \Gamma \fCenter \Delta', !B \lif !C,
!C$ (dengan melemahkan memakai $!B \lif !C$ di kanan). Hipotesis induksi
berlaku pada
\[
\AxiomC{}
\RightLabel{$\pi_1'$}
\Deduce$!B, \Gamma \fCenter \Delta', !B \lif !C, !C, !A$
\AxiomC{}
\RightLabel{$\pi_3'$}
\Deduce$!A, !B, \Gamma \fCenter \Delta', !B \lif !C, !C$
\RightLabel{\CutCS}
\BinaryInf$!B, \Gamma \fCenter \Delta', !B \lif !C, !C$
\DisplayProof
\]
dengan !!{formula} potongan~$!A$: !!{proof} ini mempunyai peringkat potongan
yang sama dan !!{height} potongan yang lebih rendah daripada~$\pi$. Ini
memberikan !!a{proof}~$\pi_4$ dalam \Log{G3c} tanpa \CutCS{} dari $!B,
\Gamma \fCenter \Delta', !B \lif !C, !C$. Terapkan $\RightR{\lif}$ untuk
memperoleh !!a{proof} dari $\Gamma \fCenter \Delta', !B \lif !C, !B \lif
!C$:
\[
\AxiomC{}
\RightLabel{$\pi_4$}
\Deduce$!B, \Gamma \fCenter \Delta', !B \lif !C, !C$
\RightLabel{$\RightR{\lif}$}
\UnaryInf$\Gamma \fCenter \Delta', !B \lif !C, !B \lif !C$
\DisplayProof
\]
Untuk memperoleh !!a{proof}~$\pi'$ dari $\Gamma \Sequent \Delta$, yakni
$\Gamma \fCenter \Delta', !B \lif !C$, terapkan
\olref[seq][inv]{prop:G3c-cont-adm} (dapat diterimanya kontraksi).

Sekarang andaikan $\pi$ berakhir dengan:
\[
\AxiomC{}
\RightLabel{$\pi_1$}
\Deduce$\lexists[x][!B(x)], \Gamma' \fCenter \Delta, !A$
\AxiomC{}
\RightLabel{$\pi_3$}
\Deduce$!A, !B(c), \Gamma' \fCenter \Delta$
\RightLabel{$\LeftR{\lexists}$}
\UnaryInf$!A, \lexists[x][!B(x)], \Gamma' \fCenter \Delta$
\RightSubproofLabel{$\pi_2$}
\RightLabel{\CutCS}
\BinaryInf$\Gamma' \fCenter \Delta$
\DisplayProof
\]
Kita kembali menerapkan hipotesis induksi pada
\[
\AxiomC{}
\RightLabel{$\pi_1[!B(c) \Sequent {}]$}
\Deduce$\lexists[x][!B(x)], !B(c), \Gamma' \fCenter \Delta, !A$
\AxiomC{}
\LeftLabel{$\pi_3[\lexists[x][!B(x)] \Sequent {}]$}
\Deduce$!A, \lexists[x][!B(x)], !B(c), \Gamma' \fCenter \Delta$
\RightLabel{\CutCS}
\BinaryInf$\lexists[x][!B(x)], !B(c), \Gamma' \fCenter \Delta$
\RightLabel{\LeftR{\lexists}}
\UnaryInf$\lexists[x][!B(x)], \lexists[x][!B(x)], \Gamma' \fCenter \Delta$
\DisplayProof
\]
Perhatikan bahwa syarat eigenvariabel terpenuhi, sebab syarat eigenvariabel
pada~\LeftR{\lexists} dalam~$\pi_2$ menjamin bahwa $c$ tidak muncul dalam
$\lexists[x][!B(x)]$, $\Gamma'$, ataupun~$\Delta$. Terakhir, terapkan
\olref[seq][inv]{prop:G3c-cont-adm}.


\paragraph{Kasus D:}
Setidaknya satu dari $\pi_1$ dan $\pi_2$ berakhir dengan aturan dua-premis
$R$ tempat $!A$ tidak prinsipal. Ada 6 kasus. Kita membahas kasus ketika
$!A$ tidak prinsipal dalam inferensi terakhir~$\pi_2$, yang berakhir
dengan~$\RightR{\land}$; kasus-kasus lain analog.

Andaikan $\pi$ berakhir dengan:
\[
\AxiomC{}
\RightLabel{$\pi_1$}
\Deduce$\Gamma \fCenter \Delta', !B \land !C, !A$
\AxiomC{}
\RightLabel{$\pi_3$}
\Deduce$!A, \Gamma \fCenter \Delta', !B$
\AxiomC{}
\RightLabel{$\pi_4$}
\Deduce$!A, \Gamma \fCenter \Delta', !C$
\RightLabel{$\RightR{\land}$}
\BinaryInf$!A, \Gamma \fCenter \Delta', !B \land !C$
\RightSubproofLabel{$\pi_2$}
\RightLabel{\CutCS}
\BinaryInf$\Gamma \fCenter \Delta$
\DisplayProof
\]
Hipotesis induksi berlaku pada !!{proof}s
\[
\AxiomC{}
\RightLabel{$\pi_1'$}
\Deduce$\Gamma \fCenter \Delta', !B \land !C, !B, !A$
\AxiomC{}
\RightLabel{$\pi_3'$}
\Deduce$!A, \Gamma \fCenter \Delta', !B \land !C, !B$
\RightLabel{\CutCS}
\BinaryInf$\Gamma \fCenter \Delta', !B \land !C, !B$
\DisplayProof
\]
dan
\[
\AxiomC{}
\RightLabel{$\pi_1''$}
\Deduce$\Gamma \fCenter \Delta', !B \land !C, !C, !A$
\AxiomC{}
\RightLabel{$\pi_4'$}
\Deduce$!A, \Gamma \fCenter \Delta', !B \land !C, !C$
\RightLabel{\CutCS}
\BinaryInf$\Gamma \fCenter \Delta', !B \land !C, !C$
\DisplayProof
\]
Di sini, !!{proof}s \Log{G3c} $\pi_1'$ dan $\pi_1''$ diperoleh dari~$\pi_1$
melalui pelemahan yang mempertahankan !!{height}
(\olref[seq][adm]{prop:weak-G3c-adm}), sekali dengan pelemahan di kanan
  memakai $!B$, dan sekali memakai $!C$. !!{proof}s $\pi_3'$ dan~$\pi_4'$
  diperoleh dari $\pi_3$ dan~$\pi_4$, juga melalui pelemahan yang
  mempertahankan !!{height} dari $\pi_3$ dan~$\pi_4$, masing-masing (dengan
  melemahkan memakai $!B \land !C$ di kanan).

Hipotesis induksi memberi kita !!{proof}s~$\pi_5$ dan~$\pi_6$ dalam
\Log{G3c} dari $\Gamma \fCenter \Delta', !B \land !C, !B$ dan $\Gamma
\fCenter \Delta', !B \land !C, !C$. Kita menggabungkan keduanya menjadi
!!a{proof} dari $\Gamma \Sequent \Delta', !B \land !C, !B \land !C$
dengan menggunakan aturan~$\RightR{\land}$:
\[
\AxiomC{}
\RightLabel{$\pi_5$}
\Deduce$\Gamma \fCenter \Delta', !B \land !C, !B$
\AxiomC{}
\RightLabel{$\pi_6$}
\Deduce$\Gamma \fCenter \Delta', !B \land !C, !C$
\RightLabel{\RightR{\land}}
\BinaryInf$\Gamma \fCenter \Delta', !B \land !C, !B \land !C$
\DisplayProof
\]
Untuk memperoleh !!a{proof} dari $\Gamma \Sequent \Delta$, yakni $\Gamma
\fCenter \Delta', !B \land !C$, terapkan
\olref[seq][inv]{prop:G3c-cont-adm} (dapat diterimanya kontraksi).

\paragraph{Kasus E: Mereduksi potongan.}
Kasus terakhir menyangkut keadaan ketika $\pi_1$ dan $\pi_2$ sama-sama
berakhir dengan inferensi tempat $!A$ prinsipal. Ada enam kasus, satu bagi
setiap kemungkinan !!{main operator} dari~$!A$.

Dalam setiap kasus, kita mengubah suatu \CutCS{} dengan !!{formula} potongan
$!A$ menjadi satu atau lebih potongan atas sub-!!{formula}s dari~$!A$,
yakni atas !!{formula}s aktif dari inferensi-inferensi tempat $!A$ menjadi
!!{formula} prinsipal. Karena itu, kasus-kasus ini sering disebut
``reduksi'' atau ``konversi'' inferensi potongan.

\begin{enumerate}
\item Jika $!A \ident \lnot !B$, maka $\pi$ berakhir dengan
\[
\AxiomC{}
\RightLabel{$\pi_1'$}
\Deduce$!B, \Gamma \fCenter \Delta$
\RightLabel{\RightR{\lnot}}
\UnaryInf$\Gamma \fCenter \Delta, \lnot !B$
\LeftSubproofLabel{$\pi_1$}
\AxiomC{}
\RightLabel{$\pi_2'$}
\Deduce$\Gamma \fCenter \Delta, !B$
\RightLabel{\LeftR{\lnot}}
\UnaryInf$\lnot !B, \Gamma \fCenter \Delta$
\RightSubproofLabel{$\pi_2$}
\RightLabel{\CutCS}
\BinaryInf$\Gamma \fCenter \Delta$
\DisplayProof
\]
Menurut hipotesis induksi, \CutCS{} dapat dieliminasi dari
\[
\AxiomC{}
\RightLabel{$\pi_2'$}
\Deduce$\Gamma \fCenter \Delta, !B$
\AxiomC{}
\RightLabel{$\pi_1'$}
\Deduce$!B, \Gamma \fCenter \Delta$
\RightLabel{\CutCS}
\BinaryInf$\Gamma \fCenter \Delta$
\DisplayProof
\]

\item Jika $!A \ident (!B \lor !C)$, maka $\pi$ berakhir dengan
\[
\AxiomC{}
\RightLabel{$\pi_1'$}
\Deduce$\Gamma \fCenter \Delta, !B, !C$
\RightLabel{\RightR{\lor}}
\UnaryInf$\Gamma \fCenter \Delta, !B \lor !C$
\LeftSubproofLabel{$\pi_1$}
\AxiomC{}
\RightLabel{$\pi_2'$}
\Deduce$!B, \Gamma \fCenter \Delta$
\AxiomC{}
\RightLabel{$\pi_2''$}
\Deduce$!C, \Gamma \fCenter \Delta$
\RightLabel{\LeftR{\lor}}
\BinaryInf$!B \lor !C, \Gamma \fCenter \Delta$
\RightSubproofLabel{$\pi_2$}
\RightLabel{\CutCS}
\BinaryInf$\Gamma \fCenter \Delta$
\DisplayProof
\]
Perhatikan !!{proof} yang berakhir dengan \CutCS,
\[
\AxiomC{}
\RightLabel{$\pi_1'$}
\Deduce$\Gamma \fCenter \Delta, !B, !C$
\AxiomC{}
\RightLabel{$\pi_2'''$}
\Deduce$!C, \Gamma \fCenter \Delta, !B$
\RightLabel{\CutCS}
\BinaryInf$\Gamma \fCenter \Delta, !B$
\DisplayProof
\]
dengan $\pi_2'''$ diperoleh dari $\pi_2''$ melalui pelemahan yang
mempertahankan !!{height}. Bukti itu mempunyai peringkat potongan
$< \cutrank{\pi}$ karena $\depth{!C} < \depth{!B \lor !C}$, sehingga
hipotesis induksi menghasilkan !!a{proof} $\pi_1''$ dalam \Log{G3c} dari
$\Gamma \fCenter \Delta, !B$. !!{proof} yang berakhir dengan \CutCS,
\[
\AxiomC{}
\RightLabel{$\pi_1''$}
\Deduce$\Gamma \fCenter \Delta, !B$
\AxiomC{}
\RightLabel{$\pi_2'$}
\Deduce$!B, \Gamma \fCenter \Delta$
\RightLabel{\CutCS}
\BinaryInf$\Gamma \fCenter \Delta$
\DisplayProof
\]
mempunyai peringkat potongan $= \depth{!B} < \depth{!B \lor !C}$, sehingga
hipotesis induksi kembali berlaku dan menghasilkan !!a{proof} $\pi'$ dari
$\Gamma \fCenter \Delta$.

\item $!A \ident (!B \land !C)$. Latihan.

\item Jika $!A \ident (!B \lif !C)$, maka $\pi$ berakhir dengan
\[
\AxiomC{}
\RightLabel{$\pi_1'$}
\Deduce$!B, \Gamma \fCenter \Delta, !C$
\RightLabel{\RightR{\lif}}
\UnaryInf$\Gamma \fCenter \Delta, !B \lif !C$
\LeftSubproofLabel{$\pi_1$}
\AxiomC{}
\RightLabel{$\pi_2'$}
\Deduce$\Gamma \fCenter \Delta, !B$
\AxiomC{}
\RightLabel{$\pi_2''$}
\Deduce$!C, \Gamma \fCenter \Delta$
\RightLabel{\LeftR{\lif}}
\BinaryInf$!B \lif !C, \Gamma \fCenter \Delta$
\RightSubproofLabel{$\pi_2$}
\RightLabel{\CutCS}
\BinaryInf$\Gamma \fCenter \Delta$
\DisplayProof
\]
!!{proof} yang berakhir dengan \CutCS,
\[
\AxiomC{}
\RightLabel{$\pi_1'$}
\Deduce$!B, \Gamma \fCenter \Delta, !C$
\AxiomC{}
\RightLabel{$\pi_2''$}
\Deduce$!C, \Gamma \fCenter \Delta$
\RightLabel{\CutCS}
\BinaryInf$!B, \Gamma \fCenter \Delta$
\DisplayProof
\]
mempunyai peringkat potongan yang lebih rendah daripada~$\pi$. Hipotesis
induksi menghasilkan !!a{proof}~$\pi_1''$ dari $!B, \Gamma \fCenter
\Delta$. Sekarang perhatikan !!{proof} yang berakhir dengan \CutCS,
\[
\AxiomC{}
\RightLabel{$\pi_2'$}
\Deduce$\Gamma \fCenter \Delta, !B$
\AxiomC{}
\RightLabel{$\pi_1''$}
\Deduce$!B, \Gamma \fCenter \Delta$
\RightLabel{\CutCS}
\BinaryInf$\Gamma \fCenter \Delta$
\DisplayProof
\]
Peringkat potongannya juga lebih kecil daripada peringkat potongan~$\pi$,
sehingga hipotesis induksi menghasilkan !!{proof} \Log{G3c}~$\pi'$ dari
$\Gamma \fCenter \Delta$.
\item Jika $!A \ident \lforall[x][!B(x)]$, maka $\pi$ berakhir dengan
\[
\AxiomC{}
\RightLabel{$\pi_1'$}
\Deduce$\Gamma \fCenter \Delta, !B(c)$
\RightLabel{\RightR{\lforall}}
\UnaryInf$\Gamma \fCenter \Delta, \lforall[x][!B(x)]$
\LeftSubproofLabel{$\pi_1$}
\AxiomC{}
\RightLabel{$\pi_2'$}
\Deduce$!B(t), \lforall[x][!B(x)], \Gamma \fCenter \Delta$
\RightLabel{\LeftR{\lforall}}
\UnaryInf$\lforall[x][!B(x)], \Gamma \fCenter \Delta$
\RightSubproofLabel{$\pi_2$}
\RightLabel{\CutCS}
\BinaryInf$\Gamma \fCenter \Delta$
\DisplayProof
\]
Menurut hipotesis induksi, \CutCS{} dapat dieliminasi dari
\[
\AxiomC{}
\RightLabel{$\pi_1''$}
\Deduce$!B(t), \Gamma \fCenter \Delta, \lforall[x][!B(x)]$
\AxiomC{}
\RightLabel{$\pi_2'$}
\Deduce$!B(t), \lforall[x][!B(x)], \Gamma \fCenter \Delta$
\RightLabel{\CutCS}
\BinaryInf$!B(t), \Gamma \fCenter \Delta$
\DisplayProof
\]
dengan $\pi_1''$ diperoleh dari $\pi_1$ (bukan~$\pi_1'$!) melalui pelemahan
yang mempertahankan !!{height}. (!!{proof} tersebut mempunyai peringkat
potongan yang sama, tetapi !!{height} potongan lebih rendah.) Ini memberi
kita !!a{proof}~$\pi_2''$ dari $!B(t), \Gamma \fCenter \Delta$.

!!{proof}~$\pi_1'$ mempunyai sekuen akhir $\Gamma \Sequent \Delta, !B(c)$,
dengan $c$ sebagai eigenvariabel inferensi~$\RightR{\lforall}$. Jadi, $c$
tidak muncul dalam~$!B(x)$, $\Gamma$, atau~$\Delta$. Kita mengasumsikan
!!{proof}~$\pi$ reguler, sehingga $c$ menjadi eigenvariabel hanya bagi
inferensi~$\RightR{\lforall}$ berikut ini dan hanya muncul di atasnya,
yakni hanya muncul dalam~$\pi_1'$ dan bukan eigenvariabel suatu inferensi
dalam~$\pi_1'$. Karena $c$ tidak muncul dalam~$\pi_2$, ia tidak mungkin
muncul dalam~$t$. Menurut \olref[seq][qua]{lem:subst-regular},
!!{proof}~$\Subst{\pi_1'}{t}{c}$ adalah !!a{proof} dari $\Gamma \Sequent
\Delta, !B(t)$. Sekarang kita menerapkan hipotesis induksi pada !!{proof}
\[
\AxiomC{}
\RightLabel{$\Subst{\pi_1'}{t}{c}$}
\Deduce$\Gamma \fCenter \Delta, !B(t)$
\AxiomC{}
\RightLabel{$\pi_2''$}
\Deduce$!B(t), \Gamma \fCenter \Delta$
\RightLabel{\CutCS}
\BinaryInf$\Gamma \fCenter \Delta$
\DisplayProof
\]
(Hipotesis induksi berlaku karena !!{formula} potongannya adalah~$!B(t)$,
sehingga !!{proof} tersebut mempunyai peringkat potongan lebih rendah
daripada~$\pi$.)
\item Kasus $!A \ident \lexists[x][!B(x)]$ kita serahkan sebagai latihan.
\end{enumerate}
\end{proof}

\begin{prob}
Periksa subkasus dari kasus~D dalam bukti \olref{lem:cut-adm-G3c} ketika
$!A$ tidak prinsipal dalam inferensi terakhir~$\pi_1$, yang berakhir
dengan~$\LeftR{\lif}$. (Catatan: bagian dari latihan ini ialah menyatakan
dengan jelas apa yang harus Anda buktikan, yakni bentuk $\pi$ dalam kasus
ini.)
\end{prob}

\begin{prob}
Periksa subkasus dari kasus~D dalam bukti \olref{lem:cut-adm-G3c} ketika
$!A$ tidak prinsipal dalam inferensi terakhir~$\pi_1$, yang berakhir
dengan~$\LeftR{\lforall}$.
\end{prob}

\begin{prob}
Periksa subkasus dari kasus~E dalam bukti \olref{lem:cut-adm-G3c} ketika
$!A$ prinsipal dalam inferensi terakhir $\pi_1$ maupun $\pi_2$, dan
$!A \ident (!B \land !C)$.
\end{prob}

\begin{prob}
Periksa subkasus dari kasus~E dalam bukti \olref{lem:cut-adm-G3c} ketika
$!A$ prinsipal dalam inferensi terakhir $\pi_1$ maupun $\pi_2$, dan
$!A \ident \lexists[x][!B(x)]$.
\end{prob}

Perhatikan bahwa dalam kasus~E, !!{height} potongan dari !!{proof} yang
berakhir dengan \CutCS{} dan menjadi sasaran hipotesis induksi dapat lebih
besar daripada !!{height} potongan bukti awal. Sebagai contoh, dalam kasus
$!A \ident (!B \lif !C)$, kita mengubah $\pi$ menjadi !!a{proof} yang
melibatkan dua inferensi \CutCS{},
\[
\AxiomC{}
\RightLabel{$\pi_2'$}
\Deduce$\Gamma \fCenter \Delta, !B$
\AxiomC{}
\RightLabel{$\pi_1'$}
\Deduce$!B, \Gamma \fCenter \Delta, !C$
\AxiomC{}
\RightLabel{$\pi_2''$}
\Deduce$!C, \Gamma \fCenter \Delta$
\RightLabel{\CutCS}
\BinaryInf$!B, \Gamma \fCenter \Delta$
\RightSubproofLabel{$\pi_1''$}
\RightLabel{\CutCS}
\BinaryInf$\Gamma \fCenter \Delta$
\DisplayProof
\]
Inferensi \CutCS{} atas antara $\pi_1'$ dan $\pi_2''$ mempunyai peringkat
maupun !!{height} potongan yang lebih rendah daripada~$\pi$. Namun,
!!{proof}~$\pi_1''$---hasil penerapan hipotesis induksi---tidak lagi
mempunyai !!{height} potongan yang lebih rendah daripada~$\pi$. Kendati
demikian, karena !!{formula} potongan dari \CutCS{} bawah adalah $!B$,
\emph{peringkat} potongannya lebih rendah daripada peringkat potongan~$\pi$.

\olref{lem:cut-adm-G3c} menetapkan bahwa kita dapat mengeliminasi satu
inferensi \CutCS{} teratas dari !!a{proof}. Kita dapat menunjukkan bahwa
berapa pun banyaknya inferensi \CutCS{} dapat dieliminasi dari !!a{proof}
dengan menerapkannya secara berturut-turut pada sub-!!{proof}s teratas yang
berakhir dengan~\CutCS.

\begin{thm}\ollabel{thm:cut-elim-G3c-topmost}
Setiap !!{proof} $\pi$ dalam $\Log{G3c} + \CutCS$ dapat ditransformasikan
menjadi !!a{proof} dalam~\Log{G3c}.
\end{thm}

\begin{proof}
  Melalui induksi pada banyaknya inferensi \CutCS{} dalam~$\pi$, yaitu~$n$.
  Jika $n=0$, tidak ada yang perlu dilakukan: $\pi$ sudah merupakan
  !!a{proof} dalam~\Log{G3c}. Jika tidak, perhatikan sub-!!{proof}~$\pi'$
  yang berakhir dengan inferensi \CutCS{} teratas. Subbukti itu memuat satu
  aturan \CutCS{} sebagai inferensi terakhir. Terapkan
  \olref{lem:cut-adm-G3c} untuk mentransformasikannya menjadi bukti
  bebas-\CutCS{}~$\pi''$, lalu ganti sub-!!{proof}~$\pi'$ dengan~$\pi''$.
  Hasilnya ialah !!a{proof} yang memuat $n-1$ inferensi \CutCS{}, sehingga
  hipotesis induksi berlaku.
\end{proof}

\end{document}
